% !TeX encoding = UTF-8
% !TeX spellcheck = es_ES
% !TeX root = Historia.tex
% !TEX root = Historia.tex

\subsection{Capítulo 1 Europa en la Post 2 GM}
Acabada la Segunda Guerra Mundial Europa se enfrenta a la recuperación y reconstrucción
de sus ciudades, industrias, … Europa queda dividida en dos, Occidente y Oriente. El 
tratado de rendición, y en general, para la Europa Occidental se queda más o menos igual
que antes de la guerra, pero al este de Alemania todo queda anexionado por Rusia. Siendo 
Alemania la peor parada (territorialmente) quedando dividía en la RFA (Republica Federal
Alemana, bajo el control de Francia, UK y EEUU) y la RDA (Republica Democrática Alemana,
bajo control soviético).  Así mismo la RDA queda como un satélite o provincia de Moscú
y la RFA se ve obligada a ayudar activamente a la recuperación del resto de Europa. 
 
Enseguida se vio que gracias a los ferrocarriles los Nazis pudieron desarrollar, construir
y mover, rápidamente toda su maquinaria de guerra. Por lo que se decidió potenciar este
método para mover los materiales, mano de obra y herramientas entre ciudades. En ese momento
era más barato y rápido construir una vía que arreglar una carretera. 
 
Empresarios, científicos, ingenieros y obreros alemanes fueron invitados a trabajar sin
beneficios a cambio de rebajar penas por haber colaborado con el partido nazi. Esto significó
trabajo cualificado a salarios mínimos y apertura de patentes. Así pues la tecnología creció
muy rápidamente, barata y con estructuras germanas. Así pues los modelos de tren de toda
Europa se normalizaron con los diseños que vemos ahora de la Deutsche Bahn. Aunque cada región
mantuvo su propia esencia, la mayoría de trenes que circulan en esta Europa son alemanes. 
 
Antes de la gran unificación y acabada la guerra (Época IIc -> IIIa [NEM 809]), nos encontramos
con tres tipos de empresas de ferrocarril:
\begin{itemize}
\item Las preexistentes ya antes de la guerra: Cubren distancias grandes, y suelen ser de una zona
de un país (MZA, Compañía de Trenes Andaluces, Norte, …). En un estado paupérrimo, necesitan
de mucha inversión para reconstruir tramos perdidos y otras infraestructuras.  
\item Las nuevas unificadas donde los gobiernos de cada país va creando su empresa nacional.
 Poco a poco adquiriendo los activos de las anteriores 
\item Pequeños servicios y líneas creadas ad-hoc para un proyecto de construcción y que,
además de lo que se necesite, se aprovecha para dar servicio a terceros y conseguir un poco 
de financiación extra para el proyecto. 
\end{itemize}

Muchas de estas lineas no llegaran a sobrevivir mucho tiempo despues del proyecto, algunas pocas
serán absorbidas para las redes nacionales. Pero servirán de modelo a futuras empresas, que sin
ánimo de convertirse en emporios llegan a donde la red nacional no ve ni útil, ni rentable poner
toda su infraestructura y organización. Estas empresas seguirán apareciendo y desapareciendo
durante todo el siglo XX.  Daniel Bahn comenzará más adelante como una de estas empresas.

\subsection{Capitulo 2 Comienzos Industriales}
A principios de los años 80 (Época IV - UIC) Daniel Bahn surge para dar servicio a un polígono
industrial recién terminado,  